% ===========================================================================
%  evnx for Python Backend Developers  —  Django, FastAPI, Flask, Celery
%  Persona AKS.  Build: pdflatex -interaction=nonstopmode python-backend.tex
% ===========================================================================
\documentclass[aspectratio=169,10pt]{beamer}

\newcommand{\capturedir}{../captures}
% ---------------------------------------------------------------------------
%  Shared preamble for the seven audience decks.
%
%  Each deck does, before the document body starts:
%     \newcommand{\capturedir}{../captures}
%     % ---------------------------------------------------------------------------
%  Shared preamble for the seven audience decks.
%
%  Each deck does, before the document body starts:
%     \newcommand{\capturedir}{../captures}
%     % ---------------------------------------------------------------------------
%  Shared preamble for the seven audience decks.
%
%  Each deck does, before the document body starts:
%     \newcommand{\capturedir}{../captures}
%     \input{persona-preamble}
%
%  \capturedir points at ../captures so both pools resolve:
%     \capture{24-scan.txt}              — the original v0.5.0 captures
%     \capture{persona/fe-01-scan-dist.txt} — captures recorded for these decks
% ---------------------------------------------------------------------------

\input{../common/evnx-preamble}

% ---- provenance ------------------------------------------------------------
% The original decks were pinned to tag v0.5.0 (85a52c3). Everything captured
% for the persona decks was re-run against 0.5.2, so the footline differs.
\renewcommand{\verifiedon}{\tiny\textcolor{evnxMuted}{verified against evnx 0.5.2 --- real run, not transcribed}}

% ---- audience badge on the title slide ------------------------------------
\newcommand{\audience}[1]{%
  \begin{center}\small\textcolor{evnxAccent}{\textbf{Audience: #1}}\end{center}}

% ---- verdict chips ---------------------------------------------------------
\newcommand{\solved}{\textcolor{evnxOk}{\ensuremath{\checkmark}}}
\newcommand{\partialmark}{\textcolor{evnxWarn}{\textbf{\textasciitilde}}}
\newcommand{\unsolved}{\textcolor{evnxErr}{\ensuremath{\times}}}

% ---- a correction block: something the source material got wrong -----------
\newenvironment{correction}[1]{%
  \begin{alertblock}{\small Correction --- #1}\small}{\end{alertblock}}

% ---- a trap block: works in the demo, fails in real use --------------------
\newenvironment{trap}[1]{%
  \setbeamercolor{block title}{bg=evnxWarn,fg=white}%
  \setbeamercolor{block body}{bg=evnxWarn!8}%
  \begin{block}{\small #1}\small}{\end{block}}

% ---- the persona card ------------------------------------------------------
\newcommand{\personacard}[4]{%
  \begin{block}{#1}
    \small
    \textbf{Stack:} #2 \\
    \textbf{Ships:} #3 \\
    \textbf{Loses time to:} #4
  \end{block}}

% ---- speaker note (rendered small and muted, not hidden) -------------------
\newcommand{\saythis}[1]{%
  \vspace{0.4em}{\footnotesize\textcolor{evnxMuted}{\textit{Say:} #1}}}

%
%  \capturedir points at ../captures so both pools resolve:
%     \capture{24-scan.txt}              — the original v0.5.0 captures
%     \capture{persona/fe-01-scan-dist.txt} — captures recorded for these decks
% ---------------------------------------------------------------------------

% ---------------------------------------------------------------------------
%  evnx presentation preamble — shared by all three decks
%
%  Engine: pdfLaTeX (also builds under XeLaTeX / LuaLaTeX).
%  Packages: beamer, listings, xcolor, amssymb, newunicodechar, booktabs,
%            fancyvrb, hyperref — all in texlive-latex-recommended + beamer.
%
%  \capturedir must be defined by the including document BEFORE \input of
%  this file, e.g.  \newcommand{\capturedir}{../captures}
% ---------------------------------------------------------------------------

\usetheme{Madrid}
\usecolortheme{seahorse}
\usefonttheme{professionalfonts}
\setbeamertemplate{navigation symbols}{}
\setbeamertemplate{caption}[numbered]
\setbeamerfont{frametitle}{size=\large}
\setbeamerfont{title}{size=\LARGE,series=\bfseries}

\usepackage[T1]{fontenc}
\usepackage[utf8]{inputenc}
\usepackage{lmodern}
\usepackage{amssymb}
\usepackage{booktabs}
\usepackage{xcolor}
\usepackage{listings}
\usepackage{fancyvrb}
\usepackage{newunicodechar}
\usepackage{tikz}
\usetikzlibrary{arrows.meta,positioning,shapes.geometric,fit,backgrounds}

% ---- palette ---------------------------------------------------------------
\definecolor{evnxInk}{HTML}{1F2933}
\definecolor{evnxAccent}{HTML}{2F6F4E}
\definecolor{evnxWarn}{HTML}{B7791F}
\definecolor{evnxErr}{HTML}{B03A2E}
\definecolor{evnxOk}{HTML}{2F6F4E}
\definecolor{evnxMuted}{HTML}{6B7280}
\definecolor{evnxTermBg}{HTML}{F5F5F2}
\definecolor{evnxRule}{HTML}{D1D5DB}

\setbeamercolor{structure}{fg=evnxAccent}
\setbeamercolor{alerted text}{fg=evnxErr}

% ---- glyph mapping for body text ------------------------------------------
% evnx emits exactly these 11 non-ASCII characters (enumerated from the
% captured logs, not guessed).
\newcommand{\evnxCheck}{\textcolor{evnxOk}{\ensuremath{\checkmark}}}
\newcommand{\evnxCross}{\textcolor{evnxErr}{\ensuremath{\times}}}
\newcommand{\evnxWarnSign}{\textcolor{evnxWarn}{\textbf{!}}}

\newunicodechar{✓}{\evnxCheck}
\newunicodechar{✗}{\evnxCross}
\newunicodechar{⚠}{\evnxWarnSign}
\newunicodechar{→}{\ensuremath{\rightarrow}}
\newunicodechar{↔}{\ensuremath{\leftrightarrow}}
\newunicodechar{─}{\textendash}
\newunicodechar{—}{---}
\newunicodechar{·}{\textperiodcentered}
\newunicodechar{…}{\ldots}
\newunicodechar{•}{\textbullet}
\newunicodechar{️}{}                % U+FE0F variation selector — invisible

% ---- terminal listing style ------------------------------------------------
% `literate` is required: newunicodechar does not apply inside listings.
\lstdefinestyle{evnxterm}{
  basicstyle=\ttfamily\scriptsize,
  backgroundcolor=\color{evnxTermBg},
  frame=single,
  rulecolor=\color{evnxRule},
  framesep=4pt,
  breaklines=true,
  breakatwhitespace=false,
  columns=fullflexible,
  keepspaces=true,
  showstringspaces=false,
  upquote=true,
  extendedchars=true,
  literate=
    {✓}{{\textcolor{evnxOk}{$\checkmark$}}}1
    {✗}{{\textcolor{evnxErr}{$\times$}}}1
    {⚠}{{\textcolor{evnxWarn}{\textbf{!}}}}1
    {→}{{$\rightarrow$}}1
    {↔}{{$\leftrightarrow$}}1
    {─}{{-}}1
    {—}{{---}}1
    {·}{{$\cdot$}}1
    {…}{{\ldots}}1
    {•}{{\textbullet}}1
    {️}{{}}0
    {é}{{\'e}}1
}

\lstdefinestyle{evnxcode}{
  style=evnxterm,
  basicstyle=\ttfamily\scriptsize,
  backgroundcolor=\color{white},
}

\lstset{style=evnxterm}

% ---- helper macros ---------------------------------------------------------
% \capture{file}        — include a real captured run verbatim
% \capturepart{f}{a}{b} — include lines a..b of a captured run
\newcommand{\capture}[1]{\lstinputlisting[style=evnxterm]{\capturedir/#1}}
\newcommand{\capturepart}[3]{%
  \lstinputlisting[style=evnxterm,firstline=#2,lastline=#3]{\capturedir/#1}}

% a labelled "real output" box
\newcommand{\realrun}[1]{%
  \begin{center}\scriptsize\textcolor{evnxMuted}{real run \textendash\ #1}\end{center}}

\newcommand{\cmd}[1]{\texttt{\textbf{#1}}}
\newcommand{\flag}[1]{\texttt{#1}}
\newcommand{\file}[1]{\texttt{#1}}

% exit-code chips
\newcommand{\exitzero}{\textcolor{evnxOk}{\texttt{0}}}
\newcommand{\exitone}{\textcolor{evnxWarn}{\texttt{1}}}
\newcommand{\exittwo}{\textcolor{evnxErr}{\texttt{2}}}

% a standard "verified against" footline for factual slides
\newcommand{\verifiedon}{\tiny\textcolor{evnxMuted}{verified against evnx 0.5.0 @ 9cfe75b}}

\usepackage[colorlinks=true,linkcolor=evnxAccent,urlcolor=evnxAccent]{hyperref}


% ---- provenance ------------------------------------------------------------
% The original decks were pinned to tag v0.5.0 (85a52c3). Everything captured
% for the persona decks was re-run against 0.5.2, so the footline differs.
\renewcommand{\verifiedon}{\tiny\textcolor{evnxMuted}{verified against evnx 0.5.2 --- real run, not transcribed}}

% ---- audience badge on the title slide ------------------------------------
\newcommand{\audience}[1]{%
  \begin{center}\small\textcolor{evnxAccent}{\textbf{Audience: #1}}\end{center}}

% ---- verdict chips ---------------------------------------------------------
\newcommand{\solved}{\textcolor{evnxOk}{\ensuremath{\checkmark}}}
\newcommand{\partialmark}{\textcolor{evnxWarn}{\textbf{\textasciitilde}}}
\newcommand{\unsolved}{\textcolor{evnxErr}{\ensuremath{\times}}}

% ---- a correction block: something the source material got wrong -----------
\newenvironment{correction}[1]{%
  \begin{alertblock}{\small Correction --- #1}\small}{\end{alertblock}}

% ---- a trap block: works in the demo, fails in real use --------------------
\newenvironment{trap}[1]{%
  \setbeamercolor{block title}{bg=evnxWarn,fg=white}%
  \setbeamercolor{block body}{bg=evnxWarn!8}%
  \begin{block}{\small #1}\small}{\end{block}}

% ---- the persona card ------------------------------------------------------
\newcommand{\personacard}[4]{%
  \begin{block}{#1}
    \small
    \textbf{Stack:} #2 \\
    \textbf{Ships:} #3 \\
    \textbf{Loses time to:} #4
  \end{block}}

% ---- speaker note (rendered small and muted, not hidden) -------------------
\newcommand{\saythis}[1]{%
  \vspace{0.4em}{\footnotesize\textcolor{evnxMuted}{\textit{Say:} #1}}}

%
%  \capturedir points at ../captures so both pools resolve:
%     \capture{24-scan.txt}              — the original v0.5.0 captures
%     \capture{persona/fe-01-scan-dist.txt} — captures recorded for these decks
% ---------------------------------------------------------------------------

% ---------------------------------------------------------------------------
%  evnx presentation preamble — shared by all three decks
%
%  Engine: pdfLaTeX (also builds under XeLaTeX / LuaLaTeX).
%  Packages: beamer, listings, xcolor, amssymb, newunicodechar, booktabs,
%            fancyvrb, hyperref — all in texlive-latex-recommended + beamer.
%
%  \capturedir must be defined by the including document BEFORE \input of
%  this file, e.g.  \newcommand{\capturedir}{../captures}
% ---------------------------------------------------------------------------

\usetheme{Madrid}
\usecolortheme{seahorse}
\usefonttheme{professionalfonts}
\setbeamertemplate{navigation symbols}{}
\setbeamertemplate{caption}[numbered]
\setbeamerfont{frametitle}{size=\large}
\setbeamerfont{title}{size=\LARGE,series=\bfseries}

\usepackage[T1]{fontenc}
\usepackage[utf8]{inputenc}
\usepackage{lmodern}
\usepackage{amssymb}
\usepackage{booktabs}
\usepackage{xcolor}
\usepackage{listings}
\usepackage{fancyvrb}
\usepackage{newunicodechar}
\usepackage{tikz}
\usetikzlibrary{arrows.meta,positioning,shapes.geometric,fit,backgrounds}

% ---- palette ---------------------------------------------------------------
\definecolor{evnxInk}{HTML}{1F2933}
\definecolor{evnxAccent}{HTML}{2F6F4E}
\definecolor{evnxWarn}{HTML}{B7791F}
\definecolor{evnxErr}{HTML}{B03A2E}
\definecolor{evnxOk}{HTML}{2F6F4E}
\definecolor{evnxMuted}{HTML}{6B7280}
\definecolor{evnxTermBg}{HTML}{F5F5F2}
\definecolor{evnxRule}{HTML}{D1D5DB}

\setbeamercolor{structure}{fg=evnxAccent}
\setbeamercolor{alerted text}{fg=evnxErr}

% ---- glyph mapping for body text ------------------------------------------
% evnx emits exactly these 11 non-ASCII characters (enumerated from the
% captured logs, not guessed).
\newcommand{\evnxCheck}{\textcolor{evnxOk}{\ensuremath{\checkmark}}}
\newcommand{\evnxCross}{\textcolor{evnxErr}{\ensuremath{\times}}}
\newcommand{\evnxWarnSign}{\textcolor{evnxWarn}{\textbf{!}}}

\newunicodechar{✓}{\evnxCheck}
\newunicodechar{✗}{\evnxCross}
\newunicodechar{⚠}{\evnxWarnSign}
\newunicodechar{→}{\ensuremath{\rightarrow}}
\newunicodechar{↔}{\ensuremath{\leftrightarrow}}
\newunicodechar{─}{\textendash}
\newunicodechar{—}{---}
\newunicodechar{·}{\textperiodcentered}
\newunicodechar{…}{\ldots}
\newunicodechar{•}{\textbullet}
\newunicodechar{️}{}                % U+FE0F variation selector — invisible

% ---- terminal listing style ------------------------------------------------
% `literate` is required: newunicodechar does not apply inside listings.
\lstdefinestyle{evnxterm}{
  basicstyle=\ttfamily\scriptsize,
  backgroundcolor=\color{evnxTermBg},
  frame=single,
  rulecolor=\color{evnxRule},
  framesep=4pt,
  breaklines=true,
  breakatwhitespace=false,
  columns=fullflexible,
  keepspaces=true,
  showstringspaces=false,
  upquote=true,
  extendedchars=true,
  literate=
    {✓}{{\textcolor{evnxOk}{$\checkmark$}}}1
    {✗}{{\textcolor{evnxErr}{$\times$}}}1
    {⚠}{{\textcolor{evnxWarn}{\textbf{!}}}}1
    {→}{{$\rightarrow$}}1
    {↔}{{$\leftrightarrow$}}1
    {─}{{-}}1
    {—}{{---}}1
    {·}{{$\cdot$}}1
    {…}{{\ldots}}1
    {•}{{\textbullet}}1
    {️}{{}}0
    {é}{{\'e}}1
}

\lstdefinestyle{evnxcode}{
  style=evnxterm,
  basicstyle=\ttfamily\scriptsize,
  backgroundcolor=\color{white},
}

\lstset{style=evnxterm}

% ---- helper macros ---------------------------------------------------------
% \capture{file}        — include a real captured run verbatim
% \capturepart{f}{a}{b} — include lines a..b of a captured run
\newcommand{\capture}[1]{\lstinputlisting[style=evnxterm]{\capturedir/#1}}
\newcommand{\capturepart}[3]{%
  \lstinputlisting[style=evnxterm,firstline=#2,lastline=#3]{\capturedir/#1}}

% a labelled "real output" box
\newcommand{\realrun}[1]{%
  \begin{center}\scriptsize\textcolor{evnxMuted}{real run \textendash\ #1}\end{center}}

\newcommand{\cmd}[1]{\texttt{\textbf{#1}}}
\newcommand{\flag}[1]{\texttt{#1}}
\newcommand{\file}[1]{\texttt{#1}}

% exit-code chips
\newcommand{\exitzero}{\textcolor{evnxOk}{\texttt{0}}}
\newcommand{\exitone}{\textcolor{evnxWarn}{\texttt{1}}}
\newcommand{\exittwo}{\textcolor{evnxErr}{\texttt{2}}}

% a standard "verified against" footline for factual slides
\newcommand{\verifiedon}{\tiny\textcolor{evnxMuted}{verified against evnx 0.5.0 @ 9cfe75b}}

\usepackage[colorlinks=true,linkcolor=evnxAccent,urlcolor=evnxAccent]{hyperref}


% ---- provenance ------------------------------------------------------------
% The original decks were pinned to tag v0.5.0 (85a52c3). Everything captured
% for the persona decks was re-run against 0.5.2, so the footline differs.
\renewcommand{\verifiedon}{\tiny\textcolor{evnxMuted}{verified against evnx 0.5.2 --- real run, not transcribed}}

% ---- audience badge on the title slide ------------------------------------
\newcommand{\audience}[1]{%
  \begin{center}\small\textcolor{evnxAccent}{\textbf{Audience: #1}}\end{center}}

% ---- verdict chips ---------------------------------------------------------
\newcommand{\solved}{\textcolor{evnxOk}{\ensuremath{\checkmark}}}
\newcommand{\partialmark}{\textcolor{evnxWarn}{\textbf{\textasciitilde}}}
\newcommand{\unsolved}{\textcolor{evnxErr}{\ensuremath{\times}}}

% ---- a correction block: something the source material got wrong -----------
\newenvironment{correction}[1]{%
  \begin{alertblock}{\small Correction --- #1}\small}{\end{alertblock}}

% ---- a trap block: works in the demo, fails in real use --------------------
\newenvironment{trap}[1]{%
  \setbeamercolor{block title}{bg=evnxWarn,fg=white}%
  \setbeamercolor{block body}{bg=evnxWarn!8}%
  \begin{block}{\small #1}\small}{\end{block}}

% ---- the persona card ------------------------------------------------------
\newcommand{\personacard}[4]{%
  \begin{block}{#1}
    \small
    \textbf{Stack:} #2 \\
    \textbf{Ships:} #3 \\
    \textbf{Loses time to:} #4
  \end{block}}

% ---- speaker note (rendered small and muted, not hidden) -------------------
\newcommand{\saythis}[1]{%
  \vspace{0.4em}{\footnotesize\textcolor{evnxMuted}{\textit{Say:} #1}}}


\title{evnx for Python backend developers}
\subtitle{Django, FastAPI, Celery --- and the config bugs that reach production}
\date{}

\begin{document}

\begin{frame}[plain]
  \titlepage
  \audience{Django · FastAPI · Flask · Celery · Docker Compose · Kubernetes}
  \begin{center}\footnotesize
  \textcolor{evnxMuted}{evnx 0.5.2. Every terminal block is a real run.}
  \end{center}
\end{frame}

% =========================================================================
\begin{frame}{This deck in one slide}
  \begin{block}{Python backend is evnx's strongest fit --- and here is why}
    evnx's validation rules were written from the same incident reports you
    have lived through. \texttt{DEBUG="False"}. A placeholder
    \texttt{SECRET\_KEY} in staging. A \file{.env} that Gunicorn never loaded.
    These are not generic ``file hygiene''. They are Django bugs with names.
  \end{block}

  \vspace{0.6em}
  \begin{block}{Two things the room will have been told wrong}
    \begin{enumerate}
      \item ``evnx cannot store multi-line values like PEM keys.''
            \textbf{It can.} We will run it.
      \item ``\cmd{validate} warns about \texttt{localhost} in production.''
            \textbf{It does not.} It warns about \texttt{localhost} when
            \emph{Docker} is present. Different rule, different trigger.
    \end{enumerate}
  \end{block}
\end{frame}

\begin{frame}{Contents}\tableofcontents[hideallsubsections]\end{frame}

% =========================================================================
\section{Three incidents}

\begin{frame}{You have shipped all three of these}
  \small
  \begin{block}{1 --- The debug page}
    Django error pages with full settings, publicly visible.
    \texttt{DEBUG = os.getenv("DEBUG")} and the value was the \emph{string}
    \texttt{"False"} --- non-empty, therefore truthy.
  \end{block}

  \begin{block}{2 --- The half-rotated key}
    Celery tasks failing on auth. The web container got the new
    \texttt{STRIPE\_SECRET\_KEY}; the worker container still had the old one.
    Different \texttt{env\_file:} entries in Compose.
  \end{block}

  \begin{block}{3 --- The variable nobody told DevOps about}
    \texttt{APP\_FEATURE\_FLAGS\_URL} works locally, crashes staging at startup
    on pydantic-settings validation. It was never added to Secrets Manager.
  \end{block}

  \saythis{``Show of hands for number one.'' It is always number one.}
\end{frame}

\begin{frame}{Which layer catches which}
  \small
  \begin{tabular}{@{}lll@{}}
    \toprule
    \textbf{Incident} & \textbf{Caught by} & \textbf{When} \\
    \midrule
    \texttt{DEBUG="False"}        & \cmd{evnx validate} & pre-commit, before it exists \\
    Worker/web key drift          & \cmd{evnx cloud run} & at process start, one source \\
    Missing variable in staging   & \cmd{evnx diff} / \cmd{sync -{}-check} & CI, on the PR \\
    \bottomrule
  \end{tabular}

  \vspace{1em}
  \begin{block}{The division of labour, stated once}
    \textbf{evnx guards the file. pydantic guards the process.} You need both.
    evnx cannot know your \texttt{Settings} class; pydantic cannot know that
    your teammate's \file{.env} is missing a key until the process boots.
  \end{block}
\end{frame}

% =========================================================================
\section{validate, on a real Django file}

\begin{frame}[fragile]{The input}
\begin{lstlisting}[style=evnxcode]
# .env
DEBUG=False
SECRET_KEY=your_secret_key_here
DATABASE_URL=postgresql://u:p@localhost:5432/app
CELERY_BROKER_URL=redis://localhost:6379/0
STRIPE_SECRET_KEY=changeme
SENTRY_DSN=notaurl

# .env.example also lists APP_FEATURE_FLAGS_URL
# and there is a Dockerfile in the directory
\end{lstlisting}

  \vspace{0.4em}
  \saythis{``Count the problems. I count six. Let us see how many it finds.''}
\end{frame}

\begin{frame}[fragile]{\cmd{evnx validate -{}-validate-formats}}
  \capturepart{persona/py-01-validate.txt}{1}{24}
  \realrun{part 1 of 2}
\end{frame}

\begin{frame}[fragile]{\ldots{} and the rest}
  \capturepart{persona/py-01-validate.txt}{25}{48}
  \realrun{part 2 of 2 --- exit \exitone}
\end{frame}

\begin{frame}{What it found, and what it did not}
  \small
  \begin{tabular}{@{}clp{6.0cm}@{}}
    \toprule
    & \textbf{Issue} & \textbf{Rule} \\
    \midrule
    \solved & \texttt{APP\_FEATURE\_FLAGS\_URL} missing & compared against \file{.env.example} \\
    \solved & \texttt{SECRET\_KEY} placeholder + weak   & two separate rules, both fire \\
    \solved & \texttt{STRIPE\_SECRET\_KEY=changeme}     & placeholder + weak \\
    \solved & \texttt{DEBUG="False"}                    & boolean-string trap \\
    \solved & \texttt{localhost} in two URLs            & \textbf{because a Dockerfile is present} \\
    \unsolved & \texttt{SENTRY\_DSN=notaurl}            & not checked --- see next slide \\
    \bottomrule
  \end{tabular}

  \vspace{0.7em}
  {\footnotesize 5 errors, 3 warnings, exit \exitone. Every one of them is
  \flag{-{}-fix}-able except the Docker/localhost pair, which needs a human to
  decide the service name.}
\end{frame}

% =========================================================================
\section{Two rules that are narrower than you think}

\begin{frame}[fragile]{The \texttt{localhost} rule is a \emph{Docker} rule}
  Same \file{.env}, two \texttt{localhost} URLs. The only thing that changes
  between these runs is whether a \file{Dockerfile} exists.

  \vspace{0.4em}
  \capture{persona/py-08-localhost-nodocker.txt}
  \vspace{-0.3em}
  \capture{persona/py-09-localhost-docker.txt}
  \realrun{\texttt{touch Dockerfile} is the entire difference}
\end{frame}

\begin{frame}{What that rule actually is}
  \begin{correction}{it is \texttt{check\_localhost\_docker}, not a production rule}
    Gated on Docker config being present. Fires only on keys containing
    \texttt{URL}, \texttt{HOST} or \texttt{ADDR}. There is \textbf{no}
    ``localhost in production'' rule --- a \file{.env.production} full of
    \texttt{localhost} passes clean, even with \flag{-{}-strict}.
  \end{correction}

  \vspace{0.8em}
  \begin{trap}{And note the exit code}
    Both runs exit \exitzero, and the JSON says
    \texttt{"status": "passed"}. \textbf{Warnings never fail the command}, and
    \flag{-{}-strict} does not change that --- it means ``also report variables
    present but not declared''. To gate a merge on warnings, read
    \texttt{summary.warnings} from \flag{-{}-format json}.
  \end{trap}
\end{frame}

\begin{frame}[fragile]{Gating on warnings, since the exit code will not}
  \capture{persona/py-11-validate-json-warnings.txt}
  \realrun{the same two warnings, machine-readable}

  \vspace{0.4em}
\begin{lstlisting}[style=evnxcode]
# fail the job on any warning, not just errors
test "$(evnx validate --format json | jq .summary.warnings)" -eq 0
\end{lstlisting}

  {\footnotesize Machine formats go to \textbf{stdout}; the human trailer goes
  to \textbf{stderr}. That separation is why the pipe above is safe.}
\end{frame}

\begin{frame}[fragile]{The URL rule catches shape, not spelling}
\begin{lstlisting}[style=evnxcode]
DATABASE_URL=postgress://u:p@db.internal:5432/app   # typo   -> NOT flagged
API_URL=htps://api.acme.io                          # typo   -> NOT flagged
WEBHOOK_URL=notaurl                                 # no scheme -> flagged
\end{lstlisting}

  \vspace{0.5em}
  \begin{block}{Why, from the source}
    \file{src/core/spec.rs:96} --- \texttt{is\_url} splits on
    \texttt{"://"}, accepts \textbf{any} RFC-shaped scheme, and requires a
    non-empty authority. That is deliberate: it has to accept
    \texttt{redis://}, \texttt{amqp://}, \texttt{mongodb+srv://} and
    \texttt{jdbc:postgresql://}. A scheme allowlist would break more than it
    catches.
  \end{block}

  \vspace{0.4em}
  {\footnotesize \textbf{Say the honest version:} it catches ``this is not a
  URL at all'', not ``you typed \texttt{postgress}''. For the typo, use
  \cmd{evnx spec} with \texttt{format = "url"} plus your own regex.}
\end{frame}

\begin{frame}{Also: the format check is keyed on the \emph{name}}
  \texttt{SENTRY\_DSN=notaurl} was not flagged. \texttt{WEBHOOK\_URL=notaurl}
  was.

  \vspace{0.8em}
  \begin{block}{The trigger}
    Format validation applies to keys containing \texttt{URL}, \texttt{URI} or
    \texttt{ENDPOINT}. \texttt{DSN} is not in that set --- so Sentry,
    Bugsnag and anything else using ``DSN'' as a noun is invisible to it.
  \end{block}

  \vspace{0.8em}
  \begin{block}{The fix you control}
    \cmd{evnx spec} lets you declare per-variable contracts:
    \vspace{0.3em}
    \begin{center}\ttfamily\scriptsize
    {[}vars.SENTRY\_DSN{]}\\ format = "url"\\ required = true
    \end{center}
  \end{block}
\end{frame}

% =========================================================================
\section{Multi-line values}

\begin{frame}{The claim you have probably read}
  \begin{center}\Large
  ``Multi-line values (PEM keys, JSON service accounts) --- \\
  \textbf{not supported by the parser}.''
  \end{center}

  \vspace{1.2em}
  This appears in the persona research, in the backlog, and it shaped a
  roadmap item with a release window attached to it.

  \vspace{1.2em}
  \begin{center}\large It is wrong. Here is the run.\end{center}

  \saythis{This is a good moment to explain \emph{why} it was wrong: the
  verification script hit exit 2 from a missing \file{.env.example} and read it
  as a parse error.}
\end{frame}

\begin{frame}[fragile]{A PEM private key, round-tripped}
  \capture{persona/py-05-multiline-pem.txt}
  \realrun{double-quoted multi-line value}

  \vspace{0.4em}
  {\footnotesize Newlines are preserved as \texttt{\textbackslash n}, and
  \texttt{B=2} proves the parser resumed correctly on the line after the
  closing quote --- it did not swallow the rest of the file.}
\end{frame}

\begin{frame}[fragile]{A JSON service account, single-quoted}
  \capture{persona/py-07-multiline-json-sq.txt}
  \realrun{single-quoted multi-line value}

  \vspace{0.4em}
  {\footnotesize \texttt{GOOGLE\_APPLICATION\_CREDENTIALS} as an inline value
  works, provided you wrap it in \textbf{single} quotes.}
\end{frame}

\begin{frame}[fragile]{The one shape that does break}
  \capture{persona/py-06-multiline-json-dq.txt}
  \realrun{double-quoted, containing escaped \texttt{\textbackslash "}}

  \vspace{0.4em}
  \begin{trap}{The rule to remember}
    A \textbf{double}-quoted multi-line value ends at the first inner
    \texttt{"}, escaped or not. So JSON pasted inside double quotes fails.
    \textbf{Wrap JSON in single quotes.} PEM has no inner quotes, so either
    works.
  \end{trap}

  \vspace{0.3em}
  {\footnotesize Exit \exittwo\ with a line number --- it fails loudly, which
  is the right failure. Worth filing as a parser bug.}
\end{frame}

% =========================================================================
\section{Every process, one source}

\begin{frame}[fragile]{\cmd{cloud run} --- the Gunicorn/Celery family of bugs}
\begin{lstlisting}[style=evnxcode]
evnx cloud run -- gunicorn config.wsgi
evnx cloud run -- uvicorn app.main:app --reload
evnx cloud run -- celery -A app worker -l info
evnx cloud run -- pytest
\end{lstlisting}

  \vspace{0.6em}
  Web, worker, beat and tests all read the \textbf{same vault version}. The
  ``worker has the old Stripe key'' incident cannot happen, because there is no
  second file to fall behind.

  \vspace{0.8em}
  \begin{block}{Verified properties}
    The secret is \textbf{absent from} \file{/proc/self/cmdline}, present in the
    child's environment, and the child's exit code is returned unchanged ---
    so it composes with \texttt{pytest} in CI.
  \end{block}

  \vspace{0.4em}
  {\footnotesize For \emph{production}, use platform injection (Kubernetes),
  not \cmd{cloud run}. It removes the plaintext file, not the master-password
  prompt.}
\end{frame}

\begin{frame}[fragile]{\flag{-{}-include} narrows what the child sees}
\begin{lstlisting}[style=evnxcode]
# collectstatic needs config, not credentials
evnx cloud run --include "APP_*" --include "STATIC_*" \
  -- python manage.py collectstatic --noinput
\end{lstlisting}

  \vspace{0.6em}
  \begin{block}{A detail worth trusting}
    A filter that matches \textbf{nothing} is an \textbf{error}, not a silent
    run with zero secrets. That distinction is the difference between a build
    step that fails and a build step that quietly ships a broken image.
  \end{block}
\end{frame}

% =========================================================================
\section{Handing off to DevOps}

\begin{frame}[fragile]{\cmd{convert -{}-to kubernetes}}
  \capture{persona/py-04-k8s.txt}
  \realrun{\texttt{stringData:} by default; \flag{-{}-base64} gives \texttt{data:}}

  \vspace{0.3em}
  {\footnotesize \textbf{Pipe it, never commit it.} This is plaintext. Feed it
  to Sealed Secrets or External Secrets Operator.}
\end{frame}

\begin{frame}[fragile]{\cmd{migrate} generates commands --- it does not upload}
  \capturepart{persona/ops-08-migrate-live-aws.txt}{1}{20}
  \realrun{\texttt{evnx migrate -{}-to aws-secrets-manager} --- \emph{without} \flag{-{}-dry-run}}

  \vspace{0.3em}
  \begin{correction}{read the banner}
    ``Generate destination commands from your \file{.env}''. Eight of nine
    destinations \textbf{print the platform's own CLI commands and contact
    nothing}. Only \texttt{github-actions} actually transfers. This is
    deliberate --- operators review the commands before running them --- but
    it surprised people when the summary used to say ``uploaded''.
  \end{correction}
\end{frame}

\begin{frame}{Where the DevOps handoff is still manual}
  \small
  \begin{tabular}{@{}clp{6.4cm}@{}}
    \toprule
    & \textbf{Gap} & \textbf{Today's workaround} \\
    \midrule
    \partialmark & No ``compare vault vs AWS SM'' & export from AWS, \cmd{evnx diff} against the file \\
    \partialmark & Web and worker Compose wiring  & \cmd{convert -{}-to docker-compose} for one shared block \\
    \partialmark & pydantic \texttt{env\_prefix} awareness & \cmd{convert -{}-to json -{}-include "APP\_*"} to see what the class will receive \\
    \unsolved & CI without a master password & machine identities are not built \\
    \unsolved & Secrets printed in logs      & app concern --- \texttt{SecretStr}, logging filters \\
    \bottomrule
  \end{tabular}

  \vspace{0.7em}
  \saythis{``The CI one is the real blocker for adoption, and it is the number
  two item on the roadmap. I will not pretend otherwise.''}
\end{frame}

% =========================================================================
\section{Adoption}

\begin{frame}[fragile]{Copy-paste: per service}
\begin{lstlisting}[style=evnxcode]
evnx init --blueprint django          # or fastapi, python
evnx vault create billing-api --env development
evnx cloud link billing-api/development && evnx cloud push

# local: every process from one source
evnx cloud run -- uvicorn app.main:app --reload
evnx cloud run -- celery -A app worker -l info

# quality gates (pre-commit + CI)
evnx validate --strict --validate-formats
evnx scan .
evnx sync --check

# hand-off
evnx migrate --to aws-secrets-manager \
  --secret-name prod/billing-api --dry-run
evnx convert --to kubernetes > /tmp/secret.yaml   # review, never commit
\end{lstlisting}
\end{frame}

\begin{frame}[fragile]{Pair it with typed settings --- both layers}
\begin{lstlisting}[style=evnxcode,language=Python]
class Settings(BaseSettings):
    model_config = SettingsConfigDict(
        env_prefix="APP_", env_nested_delimiter="__")

    debug: bool = False              # pydantic parses "false"
    database_url: PostgresDsn        # pydantic catches postgress://
    stripe_key: SecretStr            # pydantic keeps it out of repr()
\end{lstlisting}

  \vspace{0.8em}
  \begin{block}{Note what each layer covers}
    \texttt{PostgresDsn} catches the scheme typo that \cmd{validate} lets
    through. \texttt{SecretStr} covers the logging gap evnx marks
    \unsolved. \textbf{The gaps in one are the strengths of the other} ---
    which is the argument for adopting both.
  \end{block}
\end{frame}

\begin{frame}{The scorecard, honestly}
  \begin{center}
  \begin{tabular}{@{}clp{6.2cm}@{}}
    \toprule
    & \textbf{Count} & \textbf{Items} \\
    \midrule
    \solved     & 14 & onboarding, blueprints, drift, CSPRNG secrets, boolean trap, URL shape, Docker/localhost, \cmd{cloud run}, test env, code scanning, weak secrets, gitignore, K8s, Compose \\
    \partialmark & 6 & web/worker wiring, pydantic prefixes, rotation, AWS comparison, CI credentials, build-only env \\
    \unsolved   & 1 & secrets in logs \\
    \bottomrule
  \end{tabular}
  \end{center}

  \vspace{0.6em}
  \begin{block}{Revised from the source material}
    Multi-line values were listed \unsolved. They work, with one quoting
    caveat --- so the count moved from 2 unsolved to 1.
  \end{block}
\end{frame}

\begin{frame}{What is coming, in priority order}
  \small
  \begin{tabular}{@{}clp{5.6cm}@{}}
    \toprule
    & \textbf{Feature} & \textbf{Fixes} \\
    \midrule
    1 & Schema import --- read a pydantic \texttt{BaseSettings} and validate required keys and types from it & the staging-crash incident \\
    2 & Escaped quotes in double-quoted multi-line values & the JSON parse failure \\
    3 & Shared / inherited vaults (\texttt{base} + service) & three services, one DB URL \\
    4 & \cmd{diff -{}-against aws:prod/billing-api} & vault vs Secrets Manager \\
    5 & Machine identity for CI (no master password) & pipeline adoption \\
    \bottomrule
  \end{tabular}
\end{frame}

\begin{frame}[plain]
  \vfill
  \begin{center}
    {\LARGE\textbf{Questions}}\\[1.2em]
    {\small \url{https://www.evnx.dev/guides} · \url{https://github.com/urwithajit9/evnx}}\\[0.8em]
    {\footnotesize\textcolor{evnxMuted}{evnx 0.5.2 · every terminal block is a real run}}
  \end{center}
  \vfill
\end{frame}

\end{document}
